% Course: 77972 -- Physical Models of Biological Systems
% Academic year: 2026--2027 | Semester: A
% Lecture number: 1 | Document chapter: 0
% Lecture date: 2026-10-13
% Title: מבוא למידול פיזיקלי של מערכות ביולוגיות
% Lecturer: Dr. Amir Erez
% Notes prepared by: Eran Rehani
% Revision dates: 2026-10-25; 2026-11-02 (pitfall remark; parity with the English edition).
% Editorial revision: 2026-11-05 -- Hebrew prose and retrieval practice.

\chapter{מבוא למידול פיזיקלי של מערכות ביולוגיות}\label{chap:mavo}
\noindent\textbf{מתוך הרצאה 1 \enspace|\enspace 13 באוקטובר 2026}

\vspace{0.5em}
מידול פיזיקלי קושר בין מנגנונים ביוכימיים ותאיים לבין תצפיות כמותיות. המודל מתמקד בשאלה ביולוגית מסוימת ומנסח אותה כך שאפשר לחשב ממנה תחזיות. לשם כך מגדירים את ההנחות, הקירובים ומגבלות התוקף.

\begin{itemize}[nosep]
\item \textbf{מטרות למידה:} הפרדה בין תצפיות ניסוייות, הנחות מודל ותוצאות נגזרות; בחירת סקלת תיאור הולמת; הבנת מעגל המידול המדעי.
\end{itemize}

\section{מדוע מערכות חיות זקוקות לתיאור סלקטיבי?}\label{sec:mavo-complexity}
מערכות ביולוגיות כוללות רכיבים רבים שמתקשרים זה עם זה, מסתגלות ומשפיעות על סביבתן. לדוגמה, הופעת הפוטוסינתזה שינתה את הרכב האטמוספירה של כדור הארץ ואת התנאים שבהם התפתחו צורות חיים.

תאים ואורגניזמים שומרים על ארגון מבני ותפקודי באמצעות חילוף מתמיד של חומר ואנרגיה עם סביבתם. לכן הם מערכות פתוחות. מצב מתמיד ביולוגי (\textenglish{steady state}) אינו שיווי משקל תרמודינמי (\textenglish{thermodynamic equilibrium}): במצב מתמיד יש שטפים של חילוף חומרים ופינוי חום, אף שהגדלים המאקרוסקופיים נשארים קבועים. כך נשמר החוק השני של התרמודינמיקה, המתייחס למערכת ולסביבה יחד.

\begin{remark}
מצב מתמיד אינו מצב סטטי. רמה קבועה פירושה ששטף הכניסה מאזן את שטף היציאה, ושני השטפים יכולים להיות גדולים: האיזון עלול להסתיר החלפה מהירה, וזו הסיבה שעומס נגיפי קבוע בפרק~\ref{chap:hiv-dynamics} אינו מעיד על שחלוף איטי. שימו לב גם שקבוע קצב מסוים מופיע פעמיים: פעם בקביעת הרמה המתמדת ופעם בקצב ההרפיה אליה לאחר הפרעה.
\end{remark}

\section{מה תורם מודל מתמטי להבנה ביולוגית?}\label{sec:mavo-purpose}
בבניית מודל מגדירים \textbf{משתנים} (ריכוזים, כמויות תאים), \textbf{תהליכים} (תגובות ביוכימיות, חלוקה, תמותה, הדבקה) ו\textbf{קשרים מתמטיים} ביניהם. מן המשוואות גוזרים תחזיות ובודקים אותן מול המדידות:

\begin{itemize}
\item \textbf{עקביות פנימית (\textenglish{Consistency}):} כתיבת משוואות מחייבת להגדיר את ההשערות המילוליות. כך אפשר לבדוק אם מנגנון מוצע מתיישב עם חוקי שימור ועם סדרי הגודל הנמדדים.
\item \textbf{חיזוי והפרכה (\textenglish{Prediction \& Falsification}):} מודל מספק תחזיות כמותיות שאפשר לבדוק בניסוי. פער בין התחזית למדידה מחייב לבדוק את ההנחות, את התהליכים שנכללו ואת שגיאות המדידה.
\item \textbf{הסקה של פרמטרים חבויים (\textenglish{Inference}):} מדידות ישירות במערכות חיות הן לעיתים קרובות פולשניות, הרסניות, או בלתי אפשריות טכנולוגית. מודל מאפשר להסיק קבועי קצב פנימיים מתוך תגובת המערכת להפרעה חיצונית מבוקרת (כפי שנראה בהמשך במודל HIV ובניסויי \textenglish{pulse-chase}).
\item \textbf{אינטואיציה ועקרונות משותפים (\textenglish{Intuition}):} התמקדות בתהליכים הדומיננטיים מאפשרת לזהות מבנים משותפים למערכות שונות, כגון משוב שלילי, יציבות או מגבלות על העברת מידע.
\end{itemize}

\section{בחירת סקלת התיאור ובחינת ההנחות}\label{sec:mavo-scale}
בוחרים את רמת התיאור לפי השאלה והמדידות הזמינות: מולקולות בודדות, תאים, רקמות, אוכלוסיות או מערכות אקולוגיות.

השאלות המנחות בבחירת הסקלה:
\begin{enumerate}
\item \textbf{רציף מול בדיד:} האם אפשר לתאר את מספר התאים או החלקיקים בריכוז רציף ובמשוואות קצב, או שתנודות במספרים קטנים קובעות את התנהגות המערכת?
\item \textbf{דטרמיניסטי מול סטוכסטי:} האם ניתן להסתפק במשוואות קצב של ערכי תוחלת, או שיש צורך במשוואות מאסטר ותהליכי לידה-ודעיכה (\textenglish{birth-death processes})?
\item \textbf{מרחב מול תערובת הומוגנית:} האם צריך לתאר דיפוזיה, מפלי ריכוזים או מדורים נפרדים, או שקירוב תערובת אחידה (\textenglish{well-mixed}) מספק?
\end{enumerate}

\paragraph{הבדל בין מתאם מנבא למנגנון סיבתי.}
נתונים כמותיים יכולים לחשוף מתאמים בלי לזהות את המנגנון. לדוגמה, היקף חיפושי האינטרנט על גלידה משתנה עם עונות השנה, אך המתאם לבדו אינו מסביר את הסיבה. מודל פיזיקלי מבקש לתאר את הדינמיקה שמייצרת את התצפיות.

\begin{definition}{מעגל המידול המדעי}{modeling-cycle}
תהליך החקירה במודלים פיזיקליים הוא מעגלי:
\begin{equation}
\text{ניסוח השערה ומודל} \longrightarrow \text{פתרון מתמטי/נומרי} \longrightarrow \text{אימות מול נתונים ניסוייים} \longrightarrow \text{עדכון ותיקון}
\end{equation}
התאמה טובה בין נתונים למודל תומכת בשימוש בו בסקלה ובתנאים שנבדקו. היא אינה מוכיחה שכל פרט במודל נכון.
\end{definition}

\subsection*{סיכום ובדיקת הבנה}
מודל שימושי מתחיל בשאלה מוגדרת, בוחר משתנים ותהליכים בסקלה מתאימה, ומפיק תחזית שאפשר להשוות למדידה. התאמה לנתונים תומכת במודל בתנאים שנבדקו; היא אינה מזהה בהכרח מנגנון יחיד.
\begin{enumerate}
\item מספר התאים ברקמה קבוע במשך שבוע. האם אפשר להסיק שהתאים אינם מתחלקים? נמקו לפני קריאת התשובה.
\item איזה ניסוי יבחין בין החלפה מהירה לאיטית של רכיבי מאגר שרמתו קבועה?
\item איזה חלק בטענה הבאה הוא תצפית ואיזה הוא הנחה: ``הריכוז אחיד בשלוש נקודות מדידה, ולכן נניח ערבוב אחיד בכל המדור''?
\end{enumerate}
\paragraph{כיווני תשובה.}
מספר קבוע מגביל את ההפרש בין ייצור לפינוי, אך אינו קובע כל שטף בנפרד. סימון רכיבים ומעקב אחר החלפתם, או הפרעה מבוקרת ומדידת ההרפיה, מוסיפים מידע על השחלוף. אחידות בשלוש הנקודות היא התצפית; החלתה על המדור כולו היא הנחת מודל.

\subsection*{קריאה נוספת}
\begingroup\small
\begin{itemize}
\item ד"ר אמיר ארז, חוברת הקורס \emph{Physical Models of Living Systems} (2026), פרק 0, עמ' 9--11.
\item Philip Nelson, \emph{Physical Models of Living Systems}, פרק 2: ``Physics and Biology''.
\item J.~A. Lopez and A.~Erez, ``Mathematical Modelling and Intuition in Microbiology: A Perspective,'' \emph{Environmental Microbiology} \textbf{28}(4), e70266 (2026).
\end{itemize}
\endgroup
